\chapter{Tổng quan Tài liệu}

\section{Phân loại Tổn thương Da bằng Học sâu}

Học sâu cho phân loại tổn thương da đã thay đổi đáng kể với nghiên cứu năm 2017 của Esteva và cộng sự, trong đó huấn luyện một mạng CNN Inception-v3 trên 129.450 ảnh lâm sàng và so sánh với 21 bác sĩ da liễu được chứng nhận trên các tác vụ phân biệt nhị phân. Mạng đạt hoặc vượt hiệu suất của bác sĩ da liễu, qua đó thiết lập rằng học chuyển giao từ các tập dữ liệu huấn luyện trước quy mô lớn như ImageNet có thể bù đắp cho kích thước tương đối nhỏ của các bộ sưu tập ảnh y khoa. Một xác nhận rộng hơn theo sau: Maron và cộng sự phát hiện rằng CNN liên tục vượt trội so với bác sĩ trung vị trên 100.000 ảnh dermoscopy và 112 bác sĩ da liễu tham gia, mặc dù các bác sĩ chuyên gia cá nhân vẫn duy trì khả năng cạnh tranh.

Cả hai nghiên cứu đó đều giải quyết phân loại nhị phân. Bài toán đa lớp khó hơn, đòi hỏi phân biệt giữa bảy loại tổn thương chồng lấp về mặt hình ảnh, được chính thức hóa bởi Thử thách ISIC 2018 (Nhiệm vụ 3). Được xây dựng từ bộ sưu tập HAM10000, tập dữ liệu thử thách bao gồm: Melanocytic Nevus (NV), Melanoma (MEL), Benign Keratosis (BKL), Basal Cell Carcinoma (BCC), Actinic Keratosis (AKIEC), Vascular Lesion (VASC) và Dermatofibroma (DF). Độ khó không chỉ nằm ở mặt hình ảnh --- MEL và NV chia sẻ các đặc điểm kết cấu, MEL và BKL chia sẻ các mẫu sắc tố --- mà còn ở mặt thống kê: NV chiếm khoảng 67\% mẫu huấn luyện trong khi DF và VASC cộng lại chưa đến 3\%. Bất kỳ mô hình nào được đánh giá chỉ dựa trên accuracy tổng thể sẽ có vẻ hoạt động tốt trên tập dữ liệu này trong khi gần như hoàn toàn bỏ qua các lớp thiểu số.

\section{Sự Tiến hóa của các Kiến trúc}

Xây dựng một bộ phân loại hiệu quả cho dữ liệu kiểu ISIC đòi hỏi chọn một kiến trúc vừa chuyển giao tốt từ ImageNet bất chấp khoảng cách miền vừa có thể huấn luyện được trên ngân sách GPU sẵn có với hầu hết các nhóm nghiên cứu. Bốn kiến trúc được đánh giá trong khóa luận này, mỗi kiến trúc đại diện cho một triết lý thiết kế riêng biệt, và việc hiểu sự khác biệt của chúng giúp giải thích các kết quả thực nghiệm ở Chương~4.

\textbf{ResNet} đã giải quyết vấn đề cơ bản của việc huấn luyện mạng rất sâu: nếu không có kết nối tắt (skip connection), gradient biến mất trước khi đến các lớp đầu và quá trình học bị đình trệ. Bằng cách thêm đường tắt đồng nhất xung quanh mỗi cặp lớp tích chập, He và cộng sự cho phép gradient đi qua trực tiếp các phép biến đổi phi tuyến và đến các trọng số sớm hơn. Hiệu ứng thực tiễn ngay lập tức là các mạng 50 lớp và 152 lớp lần đầu tiên có thể huấn luyện với hội tụ ổn định. Kết hợp với batch normalisation, ResNet50 trở thành mô hình cơ sở được báo cáo phổ biến nhất trong phân loại ảnh y khoa, một phần vì 25,6 triệu tham số của nó cung cấp đủ năng lực và một phần vì hành vi của nó dưới các điều kiện huấn luyện khác nhau đã được hiểu rõ.

\textbf{DenseNet} mở rộng ý tưởng kết nối tắt đến cực điểm logic: mỗi lớp nhận các bản đồ đặc trưng được nối từ \textit{tất cả} các lớp trước trong cùng một khối dày đặc. Hệ quả cho các tập dữ liệu y khoa hạn chế là có ý nghĩa --- mỗi lớp có thể truy cập các bộ phát hiện cạnh nông và các bộ mô tả kết cấu mức trung bình mà không có những đặc trưng đó bị các tích chập tiếp theo dần ghi đè. Việc tái sử dụng này làm giảm tổng số tham số cần thiết để đạt một mức accuracy nhất định. DenseNet121 đạt hiệu suất cạnh tranh với khoảng 8 triệu tham số, khiến nó trở thành CNN hiệu quả tham số nhất trong nghiên cứu này.

\textbf{EfficientNet} tiếp cận thiết kế kiến trúc theo cách khác. Thay vì thiết kế thủ công các mô hình kết nối, Tan và Le sử dụng tìm kiếm kiến trúc thần kinh (NAS) để tìm cấu trúc khối cơ sở gọn nhẹ (B0) rồi rút ra quy tắc co giãn hỗn hợp (compound scaling) liên hệ chiều sâu, chiều rộng và độ phân giải đầu vào. Quy tắc co giãn chỉ rõ cách mở rộng đồng thời cả ba chiều khi ngân sách tính toán tăng, tránh chế độ thất bại phổ biến là chỉ co giãn một chiều với chi phí của các chiều khác. EfficientNet-B0, với 5,3 triệu tham số, đạt độ chính xác cạnh tranh với các mô hình lớn gấp nhiều lần. Đối với bối cảnh triển khai lâm sàng nơi suy luận có thể chạy trên phần cứng nhúng, hiệu quả này có ý nghĩa thực tiễn quan trọng.

\textbf{Vision Transformer} đến với phân loại ảnh thông qua Dosovitskiy và cộng sự, người chia ảnh thành các patch kích thước cố định, nhúng mỗi patch thành một vector, và xử lý chuỗi bằng cơ chế tự chú ý (self-attention). Self-attention có ưu thế cấu trúc so với tích chập: mọi patch có thể chú ý đến mọi patch khác bất kể khoảng cách không gian, điều quan trọng khi tín hiệu phân biệt trải rộng khắp ảnh thay vì cục bộ hóa. Vấn đề thực tế là quy mô dữ liệu --- ViT gốc cần huấn luyện trước trên JFT-300M, vượt xa tầm với của hầu hết các nhóm hình ảnh y khoa. DeiT thu hẹp khoảng cách đó thông qua chưng cất tri thức, nhưng \textbf{Swin Transformer} tái cấu trúc tính toán một cách cơ bản hơn: chú ý được tính trong các cửa sổ cục bộ không chồng lấp dịch chuyển giữa các lớp, khôi phục biểu diễn phân cấp đa tỷ lệ của CNN và giảm độ phức tạp từ bậc hai xuống tuyến tính theo kích thước ảnh. Swin-Tiny, với 28 triệu tham số, là mô hình lớn nhất trong nghiên cứu này.

Câu hỏi rộng hơn về việc Transformer có vốn dĩ vượt trội CNN hay không đã được Liu và cộng sự giải quyết trực tiếp, khi áp dụng các cải tiến huấn luyện thời kỳ Transformer (AdamW, lịch trình cosine, augmentation mạnh) lên một cấu trúc ResNet được hiện đại hóa và phát hiện rằng ConvNeXt thu được khớp hiệu suất Swin. Hệ quả cho khóa luận này là trực tiếp: nếu phương pháp huấn luyện chiếm ưu thế so với lựa chọn kiến trúc, thì một EfficientNet được điều chuẩn tốt trên vài nghìn ảnh dermoscopy có thể vượt trội một Swin Transformer thô --- và các thí nghiệm ở Chương~4 cho thấy chính xác kết quả này.

\section{Xử lý Mất cân bằng Lớp trong Tập dữ liệu Y tế}

Buda và cộng sự đã thực hiện một so sánh có hệ thống các giải pháp mất cân bằng lớp trên nhiều tập dữ liệu hình ảnh y khoa và phát hiện rằng các mini-batch cân bằng lớp --- đạt được thông qua oversampling hoặc lấy mẫu có trọng số --- liên tục vượt trội so với điều chỉnh trọng số hàm mất mát và undersampling. Lấy mẫu Ngẫu nhiên có Trọng số, gán mỗi mẫu huấn luyện một trọng số tỷ lệ nghịch với tần suất lớp của nó, thực hiện điều này mà không loại bỏ bất kỳ dữ liệu nào: mọi mẫu vẫn có sẵn cho huấn luyện, nhưng các mẫu lớp thiểu số được rút ra thường xuyên hơn trong mỗi mini-batch.

Thiết kế hàm mất mát cung cấp một đòn bẩy bổ sung. \textbf{Focal Loss} giảm trọng số đóng góp của các ví dụ được phân loại tự tin, đẩy tín hiệu gradient về các trường hợp khó ở ranh giới nơi các lớp hiếm và phổ biến trông giống nhau. \textbf{Label Smoothing} thay thế nhãn one-hot bằng phân phối mềm, ngăn mô hình gán xác suất 1,0 cho bất kỳ lớp nào. Müller và cộng sự đã chỉ ra thực nghiệm rằng hiệu ứng calibration này rõ rệt nhất trong các tập dữ liệu có sự mơ hồ về lớp có tính cấu trúc --- chính là sự chồng lấp MEL/NV/BKL đặc trưng cho ISIC 2018.

Tăng cường dữ liệu (data augmentation) đối phó mất cân bằng từ hướng thứ ba: tạo thêm các ví dụ huấn luyện hiệu quả. \textbf{Mixup} trộn tuyến tính các cặp ảnh và nhãn của chúng để tạo ra các mẫu ảo nằm giữa các ranh giới lớp, trong khi \textbf{CutMix} đạt hiệu ứng tương tự bằng cách dán các vùng cắt hình chữ nhật giữa các ảnh. Cả hai kỹ thuật ngăn mô hình ghi nhớ các ví dụ lớp thiểu số cụ thể, đây là một rủi ro thực sự khi Dermatofibroma chỉ có 81 ảnh huấn luyện. \textbf{Random Erasing} che các vùng input hình chữ nhật một cách ngẫu nhiên, hoạt động như dropout không gian đầu vào và cải thiện độ bền với các che khuất một phần như tóc hoặc nhiễu từ dermoscope. Shorten và Khoshgoftaar cung cấp một phân loại rộng các kỹ thuật này và liên quan.

\section{Khả năng Giải thích trong AI Y tế}

Sự chấp nhận lâm sàng của một công cụ chẩn đoán AI phụ thuộc vào nhiều hơn là độ chính xác. Một mô hình nhận diện đúng melanoma cũng phải cung cấp lập luận mà bác sĩ có thể kiểm tra --- không phải vì các tính toán nội bộ của mô hình hoàn toàn có thể diễn giải, mà vì việc gán hợp lý cho vùng ảnh đúng là một bài kiểm tra tỉnh táo có ý nghĩa. \textbf{Grad-CAM} đáp ứng yêu cầu này bằng cách tính, với một dự đoán nhất định, gradient của điểm số lớp mục tiêu đối với bản đồ đặc trưng tích chập cuối cùng và sử dụng các gradient đó để tạo ra bản đồ nhiệt không gian có trọng số theo kênh. Đối với phân loại tổn thương da cụ thể, Grad-CAM hoạt động như một bài kiểm tra tính phản chứng (falsifiability test): một mô hình tập trung activation vào nội thất và ranh giới tổn thương đang chú ý đến các đặc trưng mà bác sĩ da liễu nhận ra; một mô hình thắp sáng các vạch ruler hoặc da nền thì không, bất kể accuracy trên test set là gì.

\section{Phương pháp Kết hợp trong Hình ảnh Y khoa}

Các phương pháp kết hợp (ensemble) luôn xuất hiện trong các bài nộp đạt thứ hạng cao nhất của thử thách ISIC, và lý do thì đơn giản: các mô hình đa dạng phạm sai lầm ở những vị trí khác nhau, và việc lấy trung bình các output xác suất của chúng triệt tiêu một phần sai số đó. Ganaie và cộng sự đã khảo sát literature ensemble deep learning và phát hiện rằng biểu quyết mềm (soft voting) --- lấy trung bình các vector xác suất thô trước khi lấy argmax --- có tính cạnh tranh với các chiến lược stacking phức tạp hơn trong khi đơn giản hơn nhiều để triển khai. Đối với các ứng dụng y tế, ensemble mang theo một lợi ích bổ sung: phương sai của vector xác suất trung bình báo hiệu độ không chắc chắn của dự đoán, điều có ý nghĩa lâm sàng. Một trường hợp mà cả bốn mô hình gán xác suất cao cho melanoma đòi hỏi mức độ khẩn cấp khác so với một trường hợp mà hai mô hình dự đoán melanoma và hai mô hình dự đoán nevus.

\section{Định vị Công trình Này}

Các bài nộp đạt thứ hạng cao nhất trong thử thách ISIC 2018 gốc đạt độ chính xác đa lớp cân bằng trong khoảng 85--88\%, sử dụng các ensemble lớn của các biến thể ResNet và Inception, dữ liệu huấn luyện bên ngoài và tăng cường thời gian kiểm tra mở rộng. Công trình hiện tại thực hiện sự đánh đổi khác: bốn kiến trúc, một tập huấn luyện cố định, không có dữ liệu bên ngoài và suy luận đơn lượt --- nhưng đạt $89{,}93 \pm 0{,}66\%$ accuracy và $0{,}982 \pm 0{,}001$ Macro AUC khi trung bình qua năm hạt giống ngẫu nhiên. Các mô hình đơn lẻ có hiệu suất phù hợp với các kết quả thời kỳ thử thách; lợi ích đến từ phương pháp. Nghiên cứu loại bỏ thành phần ở Chương~4 cho thấy Mixup, CutMix và lấy mẫu có trọng số cùng nhau chiếm phần lớn sự cải thiện so với baseline không điều chuẩn, gợi ý rằng chất lượng quy trình huấn luyện là yếu tố hạn chế chính cho các hệ thống thời kỳ thử thách thay vì lựa chọn kiến trúc của họ.
